\chapter{Geometry}

\section{Geometric primitives}
	\begin{lecture}{   }
	\end{lecture}
	\begin{lecture}{}
	\[
	\begin{aligned}
		\mathrm{CW} &::= \mathrm{ClockWise} \quad \text{顺时针}\\
		\mathrm{CCW} &::= \mathrm{CounterClockWise} \quad \text{逆时针}&
	\end{aligned}
	\]	
	\end{lecture}
	\kactlimport{Point.h}
	
	\begin{lecture}{}
		\[
		\begin{array}{c}
			dis = \dfrac{|(e - s) \times (p - s)|}{|e - s|}\\
			dis \leq eps \Leftrightarrow |(e - s) \times (p - s)| \leq |e - s| \cdot eps
		\end{array}
		\]
	\end{lecture}
	\kactlimport{sideOf.h}

	\kactlimport{lineDistance.h}
	
	\begin{lecture}{求直线交点}
		\[
		1/0/-1\Leftrightarrow \text{唯一交点/无交点/重合}
		\]
		\[
		\begin{array}{c}
		L_1 : s_1 + t(e_1 - s_1)\\
		L_2 : s_2 + u(e_2 - s_2)\\
		\text{交点} X = s_1 + t(e_1 - s_1) = s_2 + u(e_2 - s_2)\\
		\end{array}
		\]
		\[
		\begin{aligned}
		\overrightarrow{s_2X} \parallel \overrightarrow{s_2e_2} 
		& \Rightarrow (X - s_2) \times (e_2 - s_2) = 0\\
		& \Rightarrow (s_1 + t(e_1 - s_1) - s_2) \times (e_2 - s_2) = 0\\
		& \Rightarrow (s_1 - s_2) \times (e_2 - s_2) 
								+ t(e_1 - s_1) \times (e_2 - s_2) = 0\\
		& \Rightarrow t = \frac{(e_2 - s_2) \times (s_1 - s_2)}
													 {(e_1 - s_1) \times (e_2 - s_2)}\\
		\end{aligned}
		\]
		\[
		\begin{aligned}
		X = & s_1 
			+ \frac{(e_2 - s_2) \times (s_1 - s_2)}
						 {(e_1 - s_1) \times (e_2 - s_2)}
			\cdot (e_1 - s_1)\\
		= & \{[(e_1 - s_1) \times (e_2 - s_2) - (e_2 - s_2) \times (s_1 - s_2)] \cdot s_1 \\
			& + [(e_2 - s_2) \times (s_1 - s_2)] \cdot e_1 \}
			\cdot \frac{1}{(e_1 - s_1) \times (e_2 - s_2)}\\
		= & \{[(e_1 - s_1) \times (e_2 - s_2) + (s_1 - s_2) \times (e_2 - s_2)] \cdot s_1 \\
			& + [(e_2 - s_2) \times (s_1 - s_2)] \cdot e_1\} 
			\cdots \frac{1}{(e_1 - s_1) \times (e_2 - s_2)}\\
		= & \frac{
			[(e_1 - s_2) \times (e_2 - s_2)] \cdot s_1 
			+ [(e_2 - s_2) \times (s_1 - s_2)] \cdot e_1}
			{(e_1 - s_1) \times (e_2 - s_2)}\\
		= & \frac{p \cdot s_1 + q \cdot e_1}{d}
		\end{aligned}
		\]
	\end{lecture}
	\kactlimport{lineIntersection.h}

	\begin{lecture}{点线距离}
	\[
	\begin{aligned}
		q_{\overleftrightarrow{}} &= 
		\operatorname{proj}_{\overleftrightarrow{se}}(p)\\
		(p-q_{\overleftrightarrow{}}) \cdot (e-s) &= 0,\\
		q_{\overleftrightarrow{}} &= s + \lambda_{\overleftrightarrow{}} (e-s),\\
		(p - s - \lambda_{\overleftrightarrow{}} (e-s))\cdot (e-s) &= 0,\\
		(p-s) \cdot (e-s) &= \lambda_{\overleftrightarrow{}}((e-s) \cdot (e-s)),\\
		\lambda_{\overleftrightarrow{}} &= \dfrac{(p-s)\cdot (e-s)}{\lVert e-s \rVert^2},\\
		\lambda_{\overline{se}}	&=
		\operatorname{clamp}\left(
		\dfrac{(p-s)\cdot(e-s)}{\lVert e-s\rVert^2}, 0, 1
		\right),\\
		q_{\overline{se}} &= s+\lambda_{\overline{se}}(e-s),\\
		\operatorname{segDist}(s,e,p) &= \lVert p-q_{\overline{se}}\rVert.
	\end{aligned}
	\]
	\end{lecture}
	\kactlimport{SegmentDistance.h}

	\begin{lecture}{点在线段上}
	\[
	\begin{aligned}
		(e-s)\times(p-s)=0
		&\Longleftrightarrow p\in \overleftrightarrow{se},\\
		(p-s)\cdot(p-e)=0
		&\Longleftrightarrow p=s \ \text{or}\ p=e,\\
		(p-s)\cdot(p-e)<0
		&\Longleftrightarrow p\in \overline{se}^{\circ}.
	\end{aligned}
	\]
	\end{lecture}
	\kactlimport{OnSegment.h}

	\begin{lecture}{线段交点}
	\[
	|res|=
	\begin{cases}
		0, & \text{无交点},\\
		1, & \text{唯一交点},\\
		2, & \text{无限交点, 返回公共线段的两个端点}.
	\end{cases}
	\]

	\[
	\begin{aligned}
		os_1 &= \overrightarrow{s_2e_2}\times\overrightarrow{s_2s_1},\\
		oe_1 &= \overrightarrow{s_2e_2}\times\overrightarrow{s_2e_1},\\
		os_2 &= \overrightarrow{s_1e_1}\times\overrightarrow{s_1s_2},\\
		oe_2 &= \overrightarrow{s_1e_1}\times\overrightarrow{s_1e_2}.
	\end{aligned}
	\]

	\[
	\begin{aligned}
		\text{唯一交点}p &= s_1+t(e_1-s_1),\\
		(e_2-s_2)\times(p-s_2)&=0.
	\end{aligned}
	\]

	\[
	\begin{aligned}
	0
	&=(e_2-s_2)\times\bigl(s_1+t(e_1-s_1)-s_2\bigr)\\
	&=(e_2-s_2)\times(s_1-s_2)
		+t(e_2-s_2)\times(e_1-s_1)\\
	&=(e_2-s_2)\times(s_1-s_2)
		+t(e_2-s_2)\times\bigl((e_1-s_2)+(s_2-s_1)\bigr)\\
	&=os_1+t(oe_1-os_1).
	\end{aligned}
	\]

	\[
	\begin{aligned}
	t&=\frac{os_1}{os_1-oe_1}.
	\end{aligned}
	\]

	\[
	\begin{aligned}
	p
	&=s_1+\frac{os_1}{os_1-oe_1}(e_1-s_1)\\
	&=\frac{s_1(os_1-oe_1)+os_1(e_1-s_1)}{os_1-oe_1}\\
	&=\frac{s_1\,oe_1-e_1\,os_1}{oe_1-os_1}.
	\end{aligned}
	\]
	\end{lecture}
	\kactlimport{SegmentIntersection.h}

	\kactlimport{LineProjectionReflection.h}
	
	\kactlimport{linearTransformation.h}
	
% 	\kactlimport{Angle.h}

% \section{Circles}
% 	\kactlimport{CircleIntersection.h}
% 	\kactlimport{CircleTangents.h}
% 	% \kactlimport{CircleLine.h}
% 	\kactlimport{CirclePolygonIntersection.h}
% 	\kactlimport{circumcircle.h}
% 	\kactlimport{MinimumEnclosingCircle.h}

% \section{Polygons}
% 	\kactlimport{InsidePolygon.h}
% 	\kactlimport{PolygonArea.h}
% 	\kactlimport{PolygonCenter.h}
% 	\kactlimport{PolygonCut.h}
% 	% \kactlimport{PolygonUnion.h}
% 	\kactlimport{ConvexHull.h}
% 	\kactlimport{HullDiameter.h}
% 	\kactlimport{PointInsideHull.h}
% 	\kactlimport{LineHullIntersection.h}

% \section{Misc. Point Set Problems}
% 	\kactlimport{ClosestPair.h}
% 	% \kactlimport{ManhattanMST.h}
% 	\kactlimport{kdTree.h}
% 	% \kactlimport{DelaunayTriangulation.h}
% 	\kactlimport{FastDelaunay.h}

% \section{3D}
% 	\kactlimport{PolyhedronVolume.h}
% 	\kactlimport{Point3D.h}
% 	\kactlimport{3dHull.h}
% 	\kactlimport{sphericalDistance.h}
